\subsection{Reflexión sobre la comparativa}

\campo{¿Qué diferencias existen entre el desarrollo nativo y el multiplataforma?}

La diferencia de fondo no es tanto técnica, sino de de dónde viene el poder. En el nativo la aplicación usa directamente los componentes del sistema; en el multiplataforma hay una capa intermedia que traduce, y de esa capa salen todas las demás diferencias:
\begin{itemize}[leftmargin=*,itemsep=0.2em,topsep=0.2em]
  \item El aspecto visual. El nativo se ve nativo porque lo es.
  \item El rendimiento y el consumo. El nativo es más eficiente, siempre.
  \item El acceso al hardware. El nativo llega a todo; el multiplataforma, a lo que alguien haya escrito un binding.
  \item El tamaño de la aplicación. El nativo pesa menos.
  \item El coste de desarrollo. Aquí el multiplataforma gana sin discusión, pero con una condición: si luego necesitas algo muy nativo en una de las dos plataformas, la ventaja se evapora.
\end{itemize}

Después de comparar las seis, mi conclusión es que la decisión no es técnica, es organizativa. La pregunta que de verdad importa es \enquote{¿qué equipo tengo y qué necesitan mis usuarios?}.

\campo{¿Qué ventajas puede ofrecer una aplicación desarrollada específicamente para Android?}

Cuatro, y las cuatro son reales:
\begin{enumerate}[leftmargin=*,itemsep=0.1em]
  \item Acceso completo a todas las APIs de Android sin intermediarios.
  \item Rendimiento y consumo óptimos.
  \item Tamaño de descarga menor.
  \item Publicación inmediata en un solo sitio.
\end{enumerate}

\campo{¿Qué beneficios puede aportar una tecnología que permita desarrollar para varios sistemas operativos?}

 Principalmente el coste, y aquí quiero ser cuidadoso con la cifra: el ahorro típico está entre el 30\,\% y el 50\,\% del coste total de desarrollo, no el 80\,\% que se repite mucho y que no se cumple. Luego la velocidad de salida al mercado, un único equipo de mantenimiento, consistencia del producto y de los valores, y actualizaciones simultáneas.

\campo{¿Qué factores debería tener en cuenta una empresa antes de elegir una tecnología?}

Si tuviera que ordenarlos por peso: dónde están nuestros usuarios y qué necesitan; qué sabe hacer el equipo y qué quiere aprender; requisitos técnicos muy específicos; presupuesto y plazo; mantenimiento a 5 años vista; riesgo de proveedor; y por último accesibilidad y regulación.

\begin{ideabox}
\textbf{\color{mint!70!black}\faLightbulb\ \ Si me pidieran un consejo concreto.}
Empezaría nativo con Kotlin y Compose si el negocio es Android, y usaría KMP para compartir la lógica cuando aparezca el iOS. Empezaría con Flutter si el equipo es pequeño y la app es visual. Y evitaría .NET MAUI salvo que ya seas una casa de C\#, porque el ciclo de soporte corto y las críticas de estabilidad no compensan en un proyecto nuevo.
\end{ideabox}